\bookchapter{Глава 18: Резервное копирование и миграция данных}

\chapterintro{У моей соседки телефон упал в пруд. Когда его достали, он уже не включался. Больше всего она жалела не о самом телефоне, а о трёх годах не сохранённых фотографий — вся жизнь ребёнка от рождения до первых шагов. Слушая её, мне стало очень горько, потому что этого можно было избежать. В этой главе мы серьёзно и основательно застрахуем ваши данные.}

\booksection{I. Сначала разберитесь: что в вашем телефоне пропадёт навсегда при потере}

У разных данных разный «статус выживания», нужно это чётко понимать:

\begin{itemize}
  \item \textbf{Хранится локально на телефоне, при потере пропадёт навсегда}: фото и видео (если не включена облачная синхронизация), заметки, SMS, история звонков, загруженные файлы.
  \item \textbf{Хранится на серверах приложений, после входа на новом телефоне восстановится}: список друзей WeChat, заказы Taobao, история транзакций в банковском приложении. Эти данные хранятся на серверах соответствующих компаний, а не в вашем телефоне.
  \item \textbf{Серая зона, требует особых действий}: история чатов WeChat (хотя серверы WeChat хранят их, вы сами не можете скачать их с сервера), учётные данные мобильного банка, ключи в приложениях двухфакторной аутентификации (типа Google Authenticator).
\end{itemize}

Так что защищать в первую очередь нужно первую и третью категории.

\booksection{II. Полное резервное копирование: всё сразу}

Полное резервное копирование — это как сделать полный «снимок» телефона: текущие настройки системы, расположение на рабочем столе, данные приложений, фото, контакты — всё упаковывается вместе.

\booksubsection{Вариант 1: Локальное копирование на компьютер (самый надёжный, рекомендуется тем, у кого есть компьютер)}

\begin{itemize}
  \item \textbf{iPhone}: подключите телефон кабелем к компьютеру, откройте iTunes (старые версии Windows) или Finder (Mac), на странице устройства выберите «Создать резервную копию сейчас». Важно: поставьте галочку «Шифровать резервную копию», иначе пароли и данные здоровья не будут сохранены.
  \item \textbf{Android}: у каждого бренда своё приложение для ПК, например, Huawei Phone Assistant, Xiaomi Assistant, Samsung Smart Switch. Подключите к компьютеру и следуйте инструкциям.
  \item Плюс локального копирования — не нужно платить за облачное пространство, а восстановление гораздо быстрее, чем из облака. Минус — нужно помнить о регулярном подключении кабеля, что большинство людей не делают.
\end{itemize}

\booksubsection{Вариант 2: Облачное резервное копирование (самое удобное, рекомендуется всем)}

\begin{itemize}
  \item \textbf{iPhone}: Настройки ваше имя iCloud Резервная копия iCloud Включить. Телефон будет автоматически создавать копию, когда подключён к питанию, Wi-Fi и заблокирован. Бесплатных 5 ГБ, скорее всего, не хватит — рекомендую расширить до 50 ГБ за 6 юаней.
  \item \textbf{Android}: у каждого бренда есть своё облачное резервное копирование. Huawei: найдите в настройках «Облачное резервное копирование», Xiaomi: «Xiaomi Cloud», OPPO/vivo — аналогично. Также есть встроенное резервное копирование Google (Настройки Система Резервное копирование) — включите и его.
  \item Главный плюс облачного копирования — «не нужно думать», всё автоматически. Главный минус — восстановление занимает много времени, а бесплатного места обычно не хватает.
\end{itemize}

\booksection{III. Стратегия выборочного резервного копирования: важное сохраните отдельно}

Полное копирование — это подстраховка, но некоторые вещи настолько важны, что заслуживают отдельной копии:

\booksubsection{Фото и видео (самый ценный цифровой актив)}

\begin{itemize}
  \item \textbf{Схема двойной страховки}: включите автоматическую синхронизацию iCloud Фото / Google Фото и одновременно раз в несколько месяцев вручную копируйте фото на компьютер или внешний жёсткий диск. Так одна копия в облаке, одна локально — взаимное резервирование.
  \item \textbf{Бесплатный вариант}: Google Фото даёт 15 ГБ бесплатно, можно настроить «сжатое качество» для экономии места. Или раз в месяц копируйте все фото на компьютер через кабель.
\end{itemize}

\booksubsection{Контакты (легко упустить, но критически важно)}

\begin{itemize}
  \item Контакты iPhone по умолчанию синхронизируются с iCloud — если синхронизация включена, дополнительное копирование не требуется. Но для надёжности можно зайти на iCloud.com и экспортировать файл vCard на компьютер.
  \item Контакты Android обычно синхронизируются с Google или облаком бренда. В настройках приложения «Контакты» экспортируйте файл .vcf — это будет полная резервная копия.
\end{itemize}

\booksubsection{История чатов WeChat (отдельная обработка, высший приоритет)}

Это нужно обсуждать отдельно, потому что в полной резервной копии история чатов WeChat часто неполная, а инструменты переноса на новый телефон вообще не справляются.

\begin{itemize}
  \item \textbf{Способ 1: Перенос на новый телефон}. WeChat Я Настройки Основные Перенос и резервное копирование чатов Перенос Перенос на телефон/планшет. Старый и новый телефон в одной сети Wi-Fi, сканируйте QR-код для передачи. Внимание: до завершения переноса не выходите из WeChat на старом телефоне.
  \item \textbf{Способ 2: Резервное копирование на компьютер}. Установите WeChat для ПК, в левом нижнем углу три полоски Перенос и резервное копирование Резервное копирование и восстановление Создать резервную копию чатов на компьютере. После смены телефона можно восстановить на новый.
  \item \textbf{Способ 3: Резервное копирование в облако}. WeChat также поддерживает резервное копирование в iCloud или облачное пространство телефона, но внимание: эта копия восстанавливается только на текущем телефоне, нельзя восстановить на другое устройство. При смене телефона используйте первые два способа.
\end{itemize}

\booksection{IV. Руководство по смене телефона: со старого на новый}

Новый телефон — это радость, но если данные не перенести, становится головной болью. Три ситуации — три подхода:

\booksubsection{Обновление в рамках одного бренда (самый простой случай)}

\begin{itemize}
  \item \textbf{iPhone новый iPhone}: при активации положите старый телефон рядом — автоматически появится «Быстрое начало». Отсканируйте узор на новом телефоне старым. Все данные, настройки, расположение приложений перенесутся один в один, включая пароли Wi-Fi. Займёт около 20–30 минут.
  \item \textbf{Android одного бренда}: используйте фирменный инструмент переноса. Huawei: «Phone Clone», Xiaomi: «Mi Mover», OPPO: «Phone Switch», vivo: «EasyShare». Шаги примерно одинаковы: установите одно и то же приложение на оба телефона, отсканируйте QR-код, выберите, что переносить, запустите.
\end{itemize}

\booksubsection{Смена бренда (Android на Android, разные бренды)}

Используйте сторонний инструмент переноса или приложение старого бренда (большинство поддерживают перенос на другие Android-телефоны). Что можно перенести: контакты, фото, SMS, историю звонков, часть приложений. Что нельзя: состояние входа в приложения (придётся залогиниться заново), историю чатов WeChat (используйте собственную функцию переноса WeChat).

\booksubsection{Кроссплатформенная смена (Android на iPhone или наоборот)}

\begin{itemize}
  \item \textbf{Android iPhone}: у Apple есть приложение «Перенос на iOS». Скачайте на Android-телефон, при активации нового iPhone выберите «Перенести данные с Android». Можно перенести контакты, фото, SMS и другие базовые данные, но данные приложений — нет. Историю чатов WeChat — через собственную функцию переноса WeChat.
  \item \textbf{iPhone Android}: у каждого бренда Android есть соответствующий инструмент миграции, но, честно говоря, опыт кроссплатформенной миграции хуже, чем в рамках одной платформы. Самое важное сделайте заранее отдельно: экспорт контактов, резервное копирование фото, перенос чатов WeChat. Остальное — постепенно.
\end{itemize}

\booksection{V. Золотое правило резервного копирования: принцип 3-2-1}

Это классический принцип в области безопасности данных, простой и запоминающийся:

\begin{itemize}
  \item \textbf{3}: важные данные должны иметь минимум три копии (один оригинал и две резервные).
  \item \textbf{2}: резервные копии на двух разных типах носителей (например, одна в облаке, одна на внешнем жёстком диске).
  \item \textbf{1}: как минимум одна копия должна храниться в другом физическом месте (защита от пожара, кражи и т.д., которые могут уничтожить все данные одновременно).
\end{itemize}

Для обычного человека схема «одна копия в облаке, одна на компьютере» уже считается очень хорошей. Не усложняйте, главное — действительно сделайте это.

\begin{tcbwarn}{Особое внимание к истории чатов WeChat}
Историю чатов WeChat невозможно перенести ни одним инструментом переноса телефона! Она не входит в покрытие полной резервной копии. Необходимо использовать встроенную функцию WeChat «Перенос чатов» или резервное копирование на компьютерную версию WeChat. Многие обнаруживают, что все чаты пропали, только после смены телефона — и тогда уже поздно, потому что старый телефон, возможно, уже сброшен до заводских настроек.

\end{tcbwarn}

\begin{tcbpractice}{Потренируйтесь: проведите полную тренировку резервного копирования}
\begin{enumerate}
  \item Проверьте, включено ли облачное резервное копирование iCloud или бренда, подтвердите время последней копии.
  \item Ручное копирование фото: подключите кабель к компьютеру, скопируйте папки DCIM и Pictures на компьютер.
  \item Резервное копирование контактов: пользователи iPhone — экспортируйте vCard на iCloud.com; пользователи Android — экспортируйте .vcf в настройках приложения «Контакты».
  \item Резервное копирование чатов WeChat: откройте WeChat Я Настройки Основные Перенос и резервное копирование чатов Резервное копирование на компьютер.
  \item Скопируйте экспортированный vcf контактов и папку с фото на внешний жёсткий диск или флешку — так вы выполните часть «одна копия офлайн в другом месте» из правила 3-2-1.
\end{enumerate}

\end{tcbpractice}
